\documentclass[11pt]{article}
\usepackage[margin=0.9in]{geometry}
\usepackage{amsmath,amssymb,amsthm,hyperref,enumitem,booktabs}
\hypersetup{colorlinks=true,linkcolor=blue,urlcolor=blue}
\setlength{\emergencystretch}{1em}
\newtheorem{theorem}{Theorem}
\newcommand{\bin}[2]{\binom{#1}{#2}}
\title{All nine four-variable quadratic monomial Hilbert series}
\author{Principia Math In Progress}
\date{22 September 2026}
\begin{document}
\maketitle
\noindent\textbf{MathDB 360643. Complete intended all-power result.}
All nine conjectured Hilbert series are proved over every field.
We also compute every Hilbert coefficient with the source's sign
convention. Two resulting negative third coefficients disprove the
broader nonnegativity conjecture, MathDB 360640. This consequence
belongs to the same theorem and receives no separate campaign credit.

\section*{Task 1: the theorem and complete coefficient table}
Let \(S=k[x_1,x_2,x_3,x_4]\), let \(G\) be a graph on
\(V=\{1,2,3,4\}\), and set
\[
 I_G=(x_1^2,x_2^2,x_3^2,x_4^2,\ x_ix_j:\{i,j\}\in E(G)).
\]
In tables, \(ij\) denotes the edge \(\{i,j\}\), and \(s\geq1\)
is an integer. We use \(\bin ab=0\) when \(0\leq a<b\).

\begin{theorem}
For each of the following nine graphs and every \(s\geq1\),
\[
 \operatorname{Hilb}(S/I_G^s;t)
 =\sum_{d=0}^{2s-1}\bin{d+3}{3}t^d
             +A_G(s)t^{2s}+B_G(s)t^{2s+1},                  \tag{1}
\]
where \(A_G,B_G\) are given in Table 1. The complete Hilbert
coefficient vectors are those in Table 2, with convention
\[
 \ell(S/I_G^s)=\sum_{j=0}^4(-1)^je_j(I_G)
                             \bin{s+3-j}{4-j}.             \tag{2}
\]
These formulas hold for all positive powers, not only eventually.
\end{theorem}

\begin{center}
\renewcommand{\arraystretch}{1.18}
\begin{tabular}{c l c c}
\toprule
Row&\(E(G)\)&\(A_G(s)\)&\(B_G(s)\)\\\midrule
1&\(12\)&\(s(s+1)(2s+3)/2\)&\(2\bin{s+2}{3}\)\\
2&\(12,13\)&\(4\bin{s+2}{3}\)&\(\bin{s+1}{2}\)\\
3&\(12,34\)&\(4\bin{s+2}{3}\)&0\\
4&\(12,13,23\)&\(s(s+1)(4s+5)/6\)&0\\
5&\(12,13,14\)&\(s(2s+1)\)&\(\bin{s+1}{2}\)\\
6&\(12,23,34\)&\(s(2s+1)\)&0\\
7&\(12,23,34,14\)&\(2s\)&0\\
8&\(12,13,14,23\)&\(s(s+1)\)&0\\
9&\(12,13,14,23,24\)&\(s\)&0\\\bottomrule
\end{tabular}
\par\smallskip Table 1. All nine source Hilbert-series formulas.
\end{center}

\begin{center}
\begin{tabular}{c rrrrr c rrrrr}
\toprule
Row&\(e_0\)&\(e_1\)&\(e_2\)&\(e_3\)&\(e_4\)&
Row&\(e_0\)&\(e_1\)&\(e_2\)&\(e_3\)&\(e_4\)\\\midrule
1&16&4&0&0&0&6&16&12&5&1&0\\
2&16&8&2&0&0&7&16&12&1&\(-2\)&0\\
3&16&8&1&0&0&8&16&12&3&0&0\\
4&16&8&0&0&0&9&16&12&1&\(-1\)&0\\
5&16&12&6&1&0&&&&&&\\\bottomrule
\end{tabular}
\par\smallskip Table 2. Complete coefficient vectors in convention (2).
\end{center}
Compared with the source, \(e_1\) in rows 2 and 3 is 8, the omitted
\(e_3\)'s in rows 5 and 6 are 1, and the signs of \(e_3\) in rows
7 and 9 are negative. The main result is the proof of every series
for every power. These arithmetic corrections alone earn no credit.

\section*{Monomial membership and the degree bound}
Write \(a_i=2b_i+\epsilon_i\), where \(b_i\geq0\) and
\(\epsilon_i\in\{0,1\}\), for a monomial \(x^a\).
Let \(O=\{i:\epsilon_i=1\}\), and let
\(U=\{i:a_i>0\}\) be its support. A product of quadratic
generators can be regarded as pairs of tokens: equal-index pairs
use squares and distinct-index pairs must be edges of \(G\).
Thus \(x^a\in I_G^s\) exactly when at least \(s\) such pairs
can be chosen within its exponent capacities.

All monomials of degree below \(2s\) survive, giving the first
sum of (1). Every monomial of degree \(2s+2\) lies in the ideal.
Indeed \(|O|\) is 0, 2 or 4. If it is 0 or 2, the squares alone
give at least \(s\) pairs. If it is 4, use any edge of \(G\)
(all nine graphs are nonempty); after removing its two tokens,
the remaining squares give \(s-1\) more pairs. Every higher-degree
monomial is a multiple of a degree-\(2s+2\) monomial, so it also
lies in \(I_G^s\).

\section*{Counting degree \(2s\)}
When \(|O|=0\), squares give an exact pairing, so there are no
survivors. Suppose \(O=\{u,v\}\). We claim that
\[
 x^a\notin I_G^s\quad\Longleftrightarrow\quad
 u,v\hbox{ lie in different components of }G[U].             \tag{3}
\]
For necessity of connectivity, an exact pairing uses only edges
inside support components; each such component must contain an
even total number of tokens. If \(u,v\) are separated, their two
components have odd totals, making a pairing impossible.
Conversely, if they are connected, choose a simple path from \(u\)
to \(v\) in \(G[U]\), using each path edge once. Every internal
vertex has a positive even exponent, hence at least two available
tokens. The endpoints have at least one. The remaining exponents
are all even and can be paired using squares. This proves (3).

Fix \(O\) and a set \(T\subseteq V\setminus O\) of the positive
even-exponent vertices. Then \(U=O\cup T\),
\(\sum b_i=s-1\), \(b_i\geq1\) on \(T\), and
\(b_i=0\) outside \(O\cup T\). Stars and bars gives
\(\bin{s}{|T|+1}\) choices; it gives zero automatically if there
are too few tokens. Let \(c_t(G)\) count those pairs \((O,T)\)
with \(|O|=2\), \(|T|=t\) and separated odd vertices. Thus
the two-odd-exponent contribution is
\[
            c_0\bin{s}{1}+c_1\bin{s}{2}+c_2\bin{s}{3}.      \tag{4}
\]
These small graph counts can be obtained without ideal computations.
For any induced graph \(H\) on \(m\) vertices, its number of
separated pairs is \(\bin m2-\sum_C\bin{|C|}{2}\), summing over
its components. Consequently \(c_0=6-|E(G)|\), \(c_1\) is the
sum of this count over the four induced triples, and \(c_2\) is
the count on \(G\) itself.

It remains to count \(|O|=4\), when \(\sum b_i=s-2\).
Rows 1, 2 and 4 have an odd-size support component, so no exact
pairing is possible. All \(\bin{s+1}{3}\) monomials survive.
Rows 3, 6, 7, 8 and 9 have a perfect matching; use that matching
and then squares, so no monomial survives.
In row 5 the graph is a star centered at 1. The three odd leaves
each require an odd, hence positive, number of incident edges in
an exact pairing. Thus at least three center tokens are necessary.
They suffice: use each star edge once, and pair the rest by squares.
The survivors therefore have \(a_1=1\), equivalently \(b_1=0\),
and there are \(\bin{s}{2}\) of them.
Let \(F_G(s)\) be this four-odd contribution. We have proved
\[
 A_G(s)=c_0\bin{s}{1}+c_1\bin{s}{2}+c_2\bin{s}{3}+F_G(s).  \tag{5}
\]

\begin{center}
\renewcommand{\arraystretch}{1.18}
\begin{tabular}{c rrr c rr}
\toprule
Row&\(c_0\)&\(c_1\)&\(c_2\)&\(F_G(s)\)&\(d_0\)&\(d_1\)\\\midrule
1&5&10&5&\(\bin{s+1}{3}\)&2&0\\
2&4&7&3&\(\bin{s+1}{3}\)&0&1\\
3&4&8&4&0&0&0\\
4&3&6&3&\(\bin{s+1}{3}\)&0&0\\
5&3&3&0&\(\bin{s}{2}\)&0&1\\
6&3&4&0&0&0&0\\
7&2&0&0&0&0&0\\
8&2&2&0&0&0&0\\
9&1&0&0&0&0&0\\\bottomrule
\end{tabular}
\par\smallskip Table 3. Counts used in (5) and (6).
\end{center}

\section*{Counting degree \(2s+1\)}
There are either one or three odd exponents. One odd exponent
leaves \(s\) square pairs, so gives no survivor. With three odd
exponents, squares initially give only \(s-1\) pairs. An edge
between two odd vertices supplies the missing pair. If the odd
set \(O\) is independent, only the fourth vertex \(v\) can
connect two odd vertices. If \(v\) has at least two neighbors
in \(O\) and \(a_v\geq2\), the length-two path through \(v\)
does supply \(s\) pairs, leaving the third odd token unpaired.

In the other cases the three odd vertices lie in three distinct
support components. A pairing leaving only one token would have
to leave a token in each of these three odd components, which is
impossible. This proves the necessity as well as the sufficiency
of the preceding criterion.

Let \(d_0\) count independent triples whose remaining vertex has
at most one neighbor in the triple, and \(d_1\) those whose remaining
vertex has at least two neighbors. In the first case all four
\(b_i\)'s are free with sum \(s-1\), giving \(\bin{s+2}{3}\)
survivors. In the second, survival forces \(b_v=0\), giving
\(\bin{s+1}{2}\). Hence
\[
 B_G(s)=d_0\bin{s+2}{3}+d_1\bin{s+1}{2}.                    \tag{6}
\]
Row 1 has two such independent triples of the first type. Rows 2
and 5 each have the single triple \(\{2,3,4\}\) of the second
type. The other rows have no independent triple. This gives the
last two columns of Table 3. Substitution into (5)--(6) gives
every entry of Table 1, completing the proof of (1).

\section*{Hilbert coefficients and the positivity counterexample}
Summing (1) gives the exact polynomial
\[
          \ell(S/I_G^s)=\bin{2s+3}{4}+A_G(s)+B_G(s).         \tag{7}
\]
The five binomial polynomials in (2) have distinct degrees, so their
coefficients are unique. Expanding (7) in that basis gives Table 2.
For example, direct polynomial expansion gives
\[
 \bin{2s+3}{4}=16\bin{s+3}{4}-12\bin{s+2}{3}+\bin{s+1}{2}.
\]
Thus row 7, the four-cycle, has
\[
 \ell(S/I_G^s)=16\bin{s+3}{4}-12\bin{s+2}{3}
                         +\bin{s+1}{2}+2s,
\]
so \(e_3(I_G)=-2\) in the source's convention. This ideal is
Artinian because it contains every square, and is generated
entirely in degree two. It therefore disproves Conjecture 1 in
the preprint, Conjecture 1.1 in the final (MathDB 360640), with
all its hypotheses satisfied. Row 9 likewise gives \(e_3=-1\).
These are consequences of the proved full series, not inferences
from the source's unproved formulas or from a sign typo alone.

\section*{Tasks 2--5: four proposed follow-ups}
\begin{enumerate}[label=\textbf{Task \arabic*:},start=2,leftmargin=*]
\item Extend the parity and support-graph classification to all
five-variable ideals generated by squares and graph edges.
\item Characterize the graphs for which all Hilbert coefficients
of the corresponding quadratic ideal are nonnegative.
\item Determine the depths of the associated graded rings for
the nine ideals in Table 1.
\item Compare their ordinary-power coefficients with those of the
integral-closure filtrations, explaining when the negative signs disappear.
\end{enumerate}
These are proposed directions, not additional claimed solutions.

\section*{Sources and verification}
R. Fr\"oberg, \emph{Hilbert coefficients of quadratic algebras},
\href{https://arxiv.org/abs/2311.02860}{arXiv:2311.02860v1},
6 November 2023, Conjecture 4, p.~6; all seven pages read.
Final: \emph{Journal of Experimental Mathematics} 1(2) (2025),
315--321, \href{https://doi.org/10.56994/JXM.001.002.007}{DOI
10.56994/JXM.001.002.007}, Conjecture 3.2, pp.~319--320;
all seven article pages and both front-matter pages read.

The source supplies the nine conjectured series. Elementary monomial
membership and stars and bars are the inputs to our all-power proof.
Negative third coefficients for general ideals were already known;
the counterexamples here satisfy the additional quadratic hypothesis.

Three supplementary exact groups check 89,856 monomials, 216 graded
counts, the graph counts, and coefficient recovery from lengths.
One combined result earns one campaign credit. Review by the
authoring agent only, without independent peer review or a proof
assistant; bounded source and claim checks are recorded in the audit.
\end{document}
